%قبل از شروع مسئله ترجیحاً چیزی ننویسید. اما اگر خواستید بنویسید، در قالب لیستی (فقرات و شمارش) بنویسید.
%برای اضافه کردن نمره، بعدی از شروع‌مسئله، [x] را اضافه کنید و به جای x  نمره مورد نظر را قرار دهید. برای زیرمسئله‌ها هم به همین صورت است.
%\شروع‌مسئله[10]
\شروع‌مسئله [25]
% صورت اصلی مسئله را در این‌جا بنویسید.
برای هر یک از مسائل زیر داده‌ساختار مناسبی پیشنهاد دهید و انتخاب خود را به طور مختصر توجیه کنید. برای این کار می‌توانید از ترکیب داده‌ساختار‌هایی که خوانده‌اید نیز استفاده کنید. در این صورت داده‌ساختار خود را شرح داده و آن را از نظر زمانی تحلیل کنید. توجه کنید که حافظهٔ داده‌ساختار شما باید حداکثر از مرتبهٔ تعداد اعضای آن باشد.

\پاسخ{پاسخ هر قسمت سؤال در ادامۀ آن آمده است. توجه کنید که پاسخ سؤال‌ها یکتا نیست.
	
}
%اگر مسئله چند قسمت داشت به این صورت عمل کنید:
\شروع{زیرمسائل}
\زیرمسئله[5] % مسئلۀ ۱آ
سینا یک کتاب فروشی دارد، در این کتاب فروشی قیمت هر دو کتاب متفاوت است و برای مشتریان نیز تنها قیمت کتاب‌ها مهم است. سینا ابتدا تمامی کتاب‌های کتاب فروشی را به ترتیب قیمت وارد داده ساختار می‌کند. سینا ممکن است به دنبال یک کتاب با قیمت خاص باشد، یک کتاب بفروشد و یا کتاب‌های جدیدی به فروشگاه اضافه کند.

\پاسخ{سینا به داده ساختاری نیاز دارد که درج، حذف و جستجو در آن در زمانی مناسب انجام شود. از آن‌جا که او از قیمت کتاب‌ها استفاده می‌کند و برای مشتریان نیز تنها قیمت اهمیت دارد، او می‌تواند از داده ساختار درخت جستجو استفاده کند و کتاب‌ها را بر اساس قیمت در این درخت قرار دهد. توجه کنیم از آن‌جا که ابتدا کتاب‌ها به ترتیب قیمت وارد می‌شوند، اگر از یک درخت جستجوی دودویی ساده استفاده کنیم باعث می‌شود ارتفاع درخت از مرتبه تعداد اعضا شود که مناسب نیست. بنابراین باید از درخت‌های دودویی جست‌وجوی متوازن مثل درخت قرمز-سیاه استفاده کنیم.
	
}

\زیرمسئله[10] % مسئلۀ ۱ب
یک استاد دانشگاه که تعداد بسیار زیادی دانشجو دارد، عملکرد کلاس خود را تنها بر اساس نمرات بیشینه و کمینه می‌سنجد. او می‌خواهد داده ساختاری داشته باشد که اضافه کردن نمره و حذف نمرات بیشینه (یا کمینه) در آن در زمان 
$\mathcal{O}(\log n)$
قابل انجام باشد، همچنین گزارش نمرات بیشینه و کمینه نیز در زمان 
$\Theta(1)$
انجام شوند.

\پاسخ{
	برای این سوال یک هرم بیشینه و یک هرم کمینه در نظر می‌گیریم و نمرات را در آن‌ها وارد می‌کنیم. در هر رأس علاوه بر نمره، اندیس همان نمره در هرم دیگر را نیز نگه می‌داریم. به صورت دقیق‌تر در هنگام درج ابتدا نمرۀ جدید را به انتهای هر دو هرم اضافه می‌کنیم و اندیس هر کدام را در رأس دیگر ذخیره می‌کنیم. سپس ادامۀ عملیات درج را انجام می‌دهیم که شامل تعداد جابه‌جایی رئوس است. در جابه‌جایی دو رأس، اندیس آن‌ها را در درخت دیگر به روزرسانی می‌کنیم.
%برای تکمیل عملیات درج ع دوبرای نمونه، برای درج نمره $x$ ابتدا آن را در هرم بیشینه درج می‌کنیم (اگر جای نمرۀ دیگری نیز تغییر کرد، اندیس متناظرش در هرم دیگر را  تغییر می‌دهیم). فرض می‌کنیم اندیس $x$ برابر با $i$ باشد. سپس $x$ را در هرم کمینه درج می‌کنیم (اگر جای نمرۀ دیگری نیز تغییر کرد، اندیس متناظرش در هرم دیگر را  تغییر می‌دهیم). در انتها اندیس این نمره در هرم کمینه را برابر با $i$ قرار داده و اندیس این نمره در هرم کمینه را نیز در هرم بیشینه می‌ریزیم.
همچنین اگر بخواهیم نمرۀ بیشینه (یا کمینه) را حدف کنیم، ابتدا اندیس متناظرش در هرم کمینه را نگه می‌داریم. سپس این نمره را در هرم بیشینه حذف می‌کنیم (در حین حذف وقتی جای نمرۀ دیگری نیز تغییر کرد، اندیس متناظرش را در هرم دیگر به روز می‌کنیم). سپس نمره را از هرم کمینه حذف می‌کنیم (با توجه به این‌که اندیس آن را ذخیره کرده‌ایم). در حین حذف نیز اندیس‌های هرم دیگر را به روز نگه می‌داریم.
%، توجه کنیم که اندیس خانه متناظر در هرم  کمینه را داریم. برای این حذف کافی است، مقدار خانۀ آخر هرم کمینه را با این خانه عوض کنیم و \متن‌لاتین{heapify} را در این اندیس اعمال کنیم، مشابه قبل اندیس‌هایی که تغییر می‌کنند را در هرم دیگر بروز می‌کنیم. حذف مینیمم نیز به شکل مشابه است. نمرات بیشینه و کمینه هم در دو ریشه هرم هستند. واضخ است که اعمال بالا در مرتبه زمانی خواسته شده انجام می‌شوند.
	
}
\زیرمسئله[10] % مسئلۀ ۱پ
یک استاد دانشگاه دیگر که او نیز تعداد بسیار زیادی دانشجو دارد، عملکرد کلاس خود را تنها بر اساس میانۀ نمرات می‌سنجد. او می‌خواهد داده ساختاری داشته باشد که اضافه کردن نمره در آن در زمان 
$\mathcal{O}(\log n)$
قابل انجام باشد، همچنین گزارش میانه نیز از 
$\Theta(1)$
باشد. توجه کنید اگر تعداد اعداد مجموعه فرد باشد، میانه برابر با عضو وسط مجموعه (در حالت مرتب شده) و اگر تعداد اعداد مجموعه فرد باشد برابر با میانگین دو عضو وسط مجموعه (در حالت مرتب شده) تعریف می‌شود.

\پاسخ{برای این سوال نیز دو هرم کمینه و بیشینه در نظر می‌گیریم. طوری عمل می‌کنیم که اگر  $n$ نمره داشته باشیم، 
	$\lfloor\frac{n}{2}\rfloor$
	نمره در هرم کمنیه و 
	$\lceil\frac{n}{2}\rceil$
	نمره در هرم بیشینه قرار داشته باشد. همچنین تمامی نمرات هرم بیشینه از هرم کمینه کمتر مساوی باشند (معادلاً ریشۀ هرم بیشینه کمتر مساوی ریشۀ هرم کمینه باشد). واضح است که اگر این ساختار را داشته باشیم، گزارش میانه در زمان $\Theta(1)$ انجام می‌شود (میانگین دو ریشه یا ریشۀ هرم بیشینه). در ابتدا هنگامی که هیچ نمره‌ای نداریم این ویژگی برقرار است. پس کافی است تنها این ویژگی را هنگام درج کردن حفظ کنیم.  فرض کنید می‌خواهیم عضو $x$ را درج کنیم. اگر $x$ بیشتر مساوی ریشه هرم بیشینه بود آن را در هرم کمینه و در غیر این صورت آن را در هرم بیشینه درج می‌کنیم. با این کار تمامی نمرات هرم بیشینه از هرم کمینه کمتر مساوی می‌مانند. تنها ممکن است رابطه تعداد اعضا برقرار نباشد. این اتفاق تنها در صورتی رخ می‌دهد که پس از درج هرم کمینه یک عضو بیشتر از هرم بیشینه داشته باشد. در این حالت کافی است ریشه هرم  کمینه را حذف و آن را در هرم بیشینه درج کنیم. با این کار هم رابطه تعداد اعضا درست می‌شود و هم نمرات هرم بیشینه از هرم کمینه کمتر مساوی می‌مانند. پس توانستیم که این ویژگی را حفظ کنیم. واضح است  که این کار حداکثر اندازه دو درج و یک حذف در هرم است. پس درج در این داده ساختار در زمان
	$\mathcal{O}(\log n)$
	انجام می‌شود.
	
}

%\زیرمسئله % مسئلۀ ۱ت
\پایان{زیرمسائل}




%%%%%%%%%%%%%%%%%%%%%%%%%%%%%%%%%%%%%%%%%%%%%%%%%%%%%
% برخی دستورات مفید برای تمرین‌ها. %
%%%%%%%%%%%%%%%%%%%%%%%%%%%%%%%%%%%%%%%%%%%%%%%%%%%%%

% فرمول‌های نمونه:
%$f(x) = \Theta(x)$
%$T(x, y) \leq \log(x) + 2^y + \binom{2n}{n}$
% {\tt code\_function}


% لیست موارد بدون شماره‌گذاری:
%\شروع{فقرات}
%    \فقره مورد اول 
%        \شروع{فقرات}
%            \فقره مورد اول داخلی 
%                \شروع{فقرات}
%                    \فقره اولین مورد خیلی داخلی 
%                    \فقره دومین مورد خیلی داخلی 
%                \پایان{فقرات}
%            \فقره مورد دوم داخلی 
%        \پایان{فقرات}
%    \فقره مورد دوم 
%\پایان{فقرات}

% لیست موارد شماره‌گذاری شده:
%\شروع{شمارش}
%    \فقره یک مورد 
%    \فقره یک مورد دیگر 
%    \فقره باز هم یک مورد دیگر
%\پایان{شمارش}

% برای نوشتن یک فرمول چند خطی که یک شماره بخوره. برای حالت‌های دیگر به سایت زیر مراجعه کنید:
% https://www.overleaf.com/learn/latex/Aligning_equations_with_amsmath
%\begin{align} 
%    \begin{split}
%        (x+y)^3 &= (x+y)^2(x+y) \\
%                &= (x^2+2xy+y^2)(x+y) \\
%                &= (x^3+2x^2y+xy^2) + (x^2y+2xy^2+y^3) \\
%                &= x^3+3x^2y+3xy^2+y^3
%    \end{split}                                 
%\end{align}

% نوشتن ماتریس و آرایه دوبعدی
%\begin{align}
%    A = 
%    \begin{bmatrix}
%        A_{11} & A_{21} \\
%        A_{21} & A_{22}
%    \end{bmatrix}
%\end{align}

% اضافه کردن تصویر با شرح. برای اضافه کردن تصویر بدون شرح، فقط دستور درج‌تصویر را استفاده کنید.
%\شروع{شکل}[h!]
%\تنظیم‌ازوسط
%  \درج‌تصویر[پهنا=0.5\پهنای‌متن]{img.jpg}
%\شرح{شرح تصویر}
%\برچسب{تصویر:۱}
%\پایان{شکل}
%\پایان{زیرمسائل}

%کد زیر را ببینید.
%%\lstinputlisting{source_filename.py}
%\begin{latin}
%\begin{lstlisting}
%# Import the modules
%import sys
%import random
%
%ans = True
%
%while ans:
%question = raw_input("Ask the magic 8 ball a question: (press enter to quit) ")
%
%answers = random.randint(1,8)
%
%if question == "":
%sys.exit()
%
%elif answers == 1:
%print "It is certain"
%
%elif answers == 2:
%print "Outlook good"
%
%elif answers == 3:
%print "You may rely on it"
%
%elif answers == 4:
%print "Ask again later"
%
%elif answers == 5:
%print "Concentrate and ask again"
%
%elif answers == 6:
%print "Reply hazy, try again"
%
%elif answers == 7:
%print "My reply is no"
%
%elif answers == 8:
%print "My sources say no"
%\end{lstlisting}
%\end{latin}
